%%
%% part2-book-components/ch16-cross-refs.tex
%%
%% Chapter 16: Cross-References — cleveref, hyperref,
%% label conventions, and all reference commands.
%%

\chapter{ارجاع هوشمند}
\label{chap:cross-refs}

\section{چرا این موضوع مهم است؟}
\label{sec:cross-refs-why}

در کتاب‌های فنی و علمی، ارجاع دقیق به عناصر
مختلف (معادله، قضیه، شکل، جدول، بخش) اهمیت
ویژه‌ای دارد. خواننده باید بتواند به‌سرعت به
مرجع مورد نظر برسد و از پیوستگی مطالب مطمئن
شود.

\lr{MatinBook} از \lr{cleveref} برای ارجاع
هوشمند استفاده می‌کند. این بسته، نوع مرجع را
به‌طور خودکار تشخیص می‌دهد و پیشوند مناسب
(قضیه، شکل، جدول) را اضافه می‌کند.

\mnote{پیش از \lr{cleveref}، برای هر نوع مرجع باید دستور جداگانه‌ای می‌نوشتیم: \lr{\textbackslash ref}، \lr{\textbackslash figref}، \lr{\textbackslash tabref} و... حالا فقط \lr{\textbackslash cref} کافی است.}

\section{دستور \cmd{cref}}
\label{sec:cross-refs-cref}

دستور \cmd{cref} پرکاربردترین دستور ارجاع در
\lr{MatinBook} است. این دستور، نوع مرجع را
خودکار تشخیص می‌دهد:

\begin{latin}
\begin{minted}{latex}
|\pc{طبق}|\cref{thm:pythagoras-ch8}|\pc{، در هر مثلث قائم‌الزاویه...}|%
\end{minted}
\end{latin}

خروجی: طبق \cref{thm:pythagoras-ch8}، در هر
مثلث قائم‌الزاویه...

\mnote{دستور \lr{\textbackslash cref} پیشوند مناسب را از فایل \lr{fa-IR.sty} می‌گیرد. مثلاً برای قضیه، «قضیه»؛ برای شکل، «شکل».}

\subsection{مراجع مختلف}
\label{sec:cross-refs-types}

جدول \cref{tab:cref-types} پیشوندهای
\cmd{cref} برای انواع مختلف مرجع را نشان
می‌دهد.

\begin{matintable}{پیشوندهای \lr{cref} برای انواع مرجع}{tab:cref-types}
\begin{tblr}{
    width = \textwidth,
    colspec = {l X[l] X[l]},
    hlines, vlines,
    colsep = 8pt,
    rowsep = 4pt,
    row{1} = {font=\bfseries},
}
نوع & برچسب & خروجی \\
معادله & \lr{eq:label} & معادله ۱.۱ \\
قضیه & \lr{thm:label} & قضیه ۱.۱ \\
لم & \lr{lem:label} & لم ۱.۱ \\
نتیجه & \lr{cor:label} & نتیجه ۱.۱ \\
گزاره & \lr{prop:label} & گزاره ۱.۱ \\
تعریف & \lr{def:label} & تعریف ۱.۱ \\
مثال & \lr{ex:label} & مثال ۱.۱ \\
نکته & \lr{rem:label} & نکته ۱.۱ \\
تمرین & \lr{exc:label} & تمرین ۱.۱ \\
شکل & \lr{fig:label} & شکل ۱.۱ \\
جدول & \lr{tab:label} & جدول ۱.۱ \\
کد & \lr{code:label} & کد ۱.۱ \\
الگوریتم & \lr{alg:label} & الگوریتم ۱.۱ \\
فصل & \lr{chap:label} & فصل ۱ \\
بخش & \lr{sec:label} & بخش ۱.۱ \\
\end{tblr}
\end{matintable}

\section{دستورات کمکی}
\label{sec:cross-refs-helpers}

\lr{MatinBook} چند دستور کمکی برای ارجاع ارائه
می‌دهد که همه به \cmd{cref} واگذار می‌شوند:

\begin{itemize}
    \item \cmd{thmref\{label\}}: معادل
    \cmd{cref\{label\}} برای قضیه.
    \item \cmd{lemref\{label\}}: برای لم.
    \item \cmd{corref\{label\}}: برای نتیجه.
    \item \cmd{propref\{label\}}: برای گزاره.
    \item \cmd{defref\{label\}}: برای تعریف.
    \item \cmd{exref\{label\}}: برای مثال.
    \item \cmd{remref\{label\}}: برای نکته.
    \item \cmd{excref\{label\}}: برای تمرین.
\end{itemize}

\mnote{این دستورات از نظر عملکردی معادل \lr{\textbackslash cref} هستند و فقط برای خوانایی بیشتر سورس کد نگه داشته شده‌اند.}

\section{ارجاع به چند مرجع}
\label{sec:cross-refs-multiple}

\subsection{ارجاع گروهی}
\label{sec:cross-refs-group}

برای ارجاع به چند مرجع به‌طور هم‌زمان، کلیدها
را با کاما جدا کنید:

\begin{latin}
\begin{minted}{latex}
|\pc{طبق}|\cref{thm:pythagoras-ch8,def:graph-ch8}|\pc{، ...}|%
\end{minted}
\end{latin}

خروجی: طبق \cref{thm:pythagoras-ch8,def:graph-ch8}، ...

\mnote{در فارسی، \lr{\textbackslash cref\{a,b\}} به‌صورت «قضیه ۱.۱ و تعریف ۱.۲» نمایش داده می‌شود.}

\subsection{ارجاع بازه‌ای}
\label{sec:cross-refs-range}

برای ارجاع به یک بازه‌ی پیوسته، از
\cmd{crefrange} استفاده کنید:

\begin{latin}
\begin{minted}{latex}
|\pc{معادلات}|\crefrange{eq:abs}{eq:inner}|\pc{...}|%
\end{minted}
\end{latin}

خروجی: معادلات \crefrange{eq:abs}{eq:inner}...

\mnote{دستور \lr{\textbackslash crefrange} فقط برای مراجعی کار می‌کند که در یک دسته هستند (مثلاً سه معادله‌ی پیوسته).}

\subsection{ارجاع با \cmd{Cref}}
\label{sec:cross-refs-Cref}

دستور \cmd{Cref} نسخه‌ی «حرف اول بزرگ» است.
در فارسی تفاوتی با \cmd{cref} ندارد، چون
فارسی حالت بزرگ و کوچک ندارد.

\begin{latin}
\begin{minted}{latex}
\Cref{thm:pythagoras-ch8}
\end{minted}
\end{latin}

\mnote{در انگلیسی، \lr{\textbackslash cref} می‌شود «theorem» و \lr{\textbackslash Cref} می‌شود «Theorem». در فارسی هر دو یکسان هستند.}

\section{نام‌گذاری برچسب‌ها}
\label{sec:cross-refs-labels}

برای اینکه برچسب‌ها معنادار و یکتا باشند،
از پیشوندهای استاندارد استفاده کنید:

\begin{matintable}{پیشوندهای استاندارد برچسب}{tab:label-prefixes}
\begin{tblr}{
    width = \textwidth,
    colspec = {l X[l] X[l]},
    hlines, vlines,
    colsep = 8pt,
    rowsep = 4pt,
    row{1} = {font=\bfseries},
}
نوع & پیشوند & مثال \\
معادله & \lr{eq:} & \lr{eq:einstein} \\
قضیه & \lr{thm:} & \lr{thm:pythagoras} \\
لم & \lr{lem:} & \lr{lem:euclid} \\
نتیجه & \lr{cor:} & \lr{cor:monotone} \\
گزاره & \lr{prop:} & \lr{prop:q-dense} \\
تعریف & \lr{def:} & \lr{def:graph} \\
مثال & \lr{ex:} & \lr{ex:integral} \\
نکته & \lr{rem:} & \lr{rem:division-zero} \\
تمرین & \lr{exc:} & \lr{exc:sum-even} \\
شکل & \lr{fig:} & \lr{fig:parabola} \\
جدول & \lr{tab:} & \lr{tab:comparison} \\
کد & \lr{code:} & \lr{code:factorial} \\
الگوریتم & \lr{alg:} & \lr{alg:binary} \\
فصل & \lr{chap:} & \lr{chap:introduction} \\
بخش & \lr{sec:} & \lr{sec:motivation} \\
زیربخش & \lr{subsec:} & \lr{subsec:history} \\
\end{tblr}
\end{matintable}

\mnote{برچسب‌های معنادار، خواندن فایل‌های \lr{tex} را آسان‌تر می‌کنند و از خطاهای ارجاع جلوگیری می‌کنند.}

\section{لینک‌های \lr{hyperref}}
\label{sec:cross-refs-hyperref}

تمام ارجاعات در \lr{MatinBook} به‌صورت
\textbf{لینک فعال} هستند. یعنی خواننده
می‌تواند روی شماره‌ی مرجع کلیک کند و به آن
برود.

\begin{latin}
\begin{minted}{latex}
\RequirePackage[bookmarks]{hyperref}

\hypersetup{
    colorlinks=true,
    linkcolor=blue,
    citecolor=blue,
    urlcolor=blue
}
\end{minted}
\end{latin}

\mnote{رنگ لینک‌ها در \lr{MatinBook} آبی است. برای تغییر، فایل \lr{mb-references.sty} را ویرایش کنید.}

\subsection{بوکمارک‌های PDF}
\label{sec:cross-refs-bookmarks}

\lr{hyperref} به‌طور خودکار بوکمارک‌های
\lr{PDF} را از فصل‌ها و بخش‌ها می‌سازد. این
بوکمارک‌ها در پنل کنار \lr{PDF} نمایش داده
می‌شوند.

\mnote{گزینه‌ی \lr{bookmarksnumbered=true} در \lr{mb-references.sty} باعث می‌شود شماره‌ها در بوکمارک‌ها هم نمایش داده شوند.}

\section{مثال‌های کاربردی}
\label{sec:cross-refs-examples}

\subsection{ارجاع در متن علمی}
\label{sec:cross-refs-example-scientific}

\begin{latin}
\begin{minted}{latex}
|\pc{طبق}|\cref{thm:fundamental}|\pc{، هر عدد طبیعی بزرگ‌تر از ۱ به عوامل اول تجزیه می‌شود.}|%
|\pc{این قضیه در}|\cref{chap:numbertheory}|\pc{به‌طور کامل بررسی شده است.}|%
|\pc{برای مثال،}|\cref{ex:primes-under-50}|\pc{اعداد اول کوچک‌تر از ۵۰ را نشان می‌دهد.}|%
\end{minted}
\end{latin}

\subsection{ارجاع در متن فنی}
\label{sec:cross-refs-example-technical}

\begin{latin}
\begin{minted}{latex}
|\pc{الگوریتم}|\cref{alg:binary}|\pc{روش جستجوی دودویی را نشان می‌دهد.}|%
|\pc{پیچیدگی این الگوریتم}|\lr{$O(\log n)$}|\pc{است، همان‌طور که در}|\cref{tab:complexity}|\pc{نشان داده شده.}|%
|\pc{پیاده‌سازی آن در}|\cref{code:binary-search}|\pc{آمده است.}|%
\end{minted}
\end{latin}

\mnote{استفاده از \lr{\textbackslash cref} در متن‌های علمی و فنی، خوانایی و دقت ارجاعات را افزایش می‌دهد.}

\section{خطاهای رایج}
\label{sec:cross-refs-mistakes}

\subsection{استفاده از \cmd{ref} به‌جای \cmd{cref}}
\label{sec:cross-refs-mistake-ref}

\textbf{مشکل:} \cmd{ref} فقط شماره‌ی مرجع را
نمایش می‌دهد، بدون پیشوند:

\begin{latin}
\begin{minted}{latex}
% |\pc{اشتباه}|:
|\pc{طبق معادله}|\ref{eq:einstein}|\pc{...}|%

% |\pc{درست}|:
|\pc{طبق}|\cref{eq:einstein}|\pc{...}|%
\end{minted}
\end{latin}

\textbf{راه‌حل:} همیشه از \cmd{cref} استفاده
کنید.

\subsection{قرار دادن \cmd{label} بیرون از محیط}
\label{sec:cross-refs-mistake-label}

\textbf{مشکل:} اگر \cmd{label} را بیرون از
محیط (مثل \env{theorem}) قرار دهید، ارجاع
کار نمی‌کند.

\textbf{راه‌حل:} \cmd{label} را همیشه داخل محیط
قرار دهید.

\subsection{برچسب تکراری}
\label{sec:cross-refs-mistake-duplicate}

\textbf{مشکل:} اگر دو عنصر برچسب یکسان داشته
باشند، \lr{LaTeX} خطا می‌دهد و ارجاع به یکی
از آن‌ها ممکن است اشتباه باشد.

\textbf{راه‌حل:} از برچسب‌های یکتا استفاده
کنید.

\subsection{ارجاع به عنصر قبل از تعریف}
\label{sec:cross-refs-mistake-forward}

\textbf{مشکل:} اگر به عنصری ارجاع دهید که
هنوز تعریف نشده، خروجی «؟؟» نمایش داده
می‌شود.

\textbf{راه‌حل:} حداقل دو بار کامپایل کنید.
اگر پس از دو بار همچنان «؟؟» بود، برچسب
را بررسی کنید.

\mnote{اگر خروجی «؟؟» دیدید، اول مطمئن شوید برچسب وجود دارد، سپس دوباره کامپایل کنید.}

\section{تمرین‌ها}
\label{sec:cross-refs-exercises}

\begin{exercise}
    به یک معادله، یک قضیه و یک شکل ارجاع
    دهید و پیشوندهای خروجی را بررسی کنید.
    \label{exc:cross-refs-basic}
\end{exercise}

\begin{exercise}
    به سه معادله به‌طور هم‌زمان ارجاع دهید
    و سپس به یک بازه‌ی پیوسته ارجاع دهید.
    \label{exc:cross-refs-multiple}
\end{exercise}

\begin{exercise}
    از دستورات کمکی (\cmd{thmref}،
    \cmd{defref} و...) استفاده کنید و خروجی
    را با \cmd{cref} مقایسه کنید.
    \label{exc:cross-refs-helpers}
\end{exercise}

\begin{exercise}
    یک برچسب تکراری بسازید و خطای آن را
    ببینید. سپس آن را اصلاح کنید.
    \label{exc:cross-refs-duplicate}
\end{exercise}

\section{خلاصه}
\label{sec:cross-refs-summary}

در این فصل، با ارجاع هوشمند در \lr{MatinBook}
آشنا شدیم:

\begin{itemize}
    \item \lr{MatinBook} از \lr{cleveref} برای
    ارجاع هوشمند استفاده می‌کند.

    \item دستور اصلی، \cmd{cref} است که نوع
    مرجع را خودکار تشخیص می‌دهد.

    \item پیشوندها از فایل \lr{fa-IR.sty}
    می‌آیند (قضیه، شکل، جدول، ...).

    \item \cmd{Cref} نسخه‌ی «حرف اول بزرگ» است
    ولی در فارسی تفاوتی با \cmd{cref} ندارد.

    \item \cmd{crefrange} برای ارجاع بازه‌ای
    و \cmd{cref\{a,b\}} برای ارجاع گروهی.

    \item دستورات کمکی: \cmd{thmref}،
    \cmd{lemref}، \cmd{corref}، \cmd{propref}،
    \cmd{defref}، \cmd{exref}، \cmd{remref}،
    \cmd{excref}.

    \item برچسب‌ها با پیشوندهای استاندارد:
    \lr{eq:}، \lr{thm:}، \lr{fig:}، \lr{tab:}،
    ...

    \item لینک‌های \lr{hyperref} آبی هستند
    و بوکمارک‌های \lr{PDF} خودکار ساخته
    می‌شوند.

    \item خطاهای رایج شامل استفاده از
    \cmd{ref} به‌جای \cmd{cref}، قرار دادن
    \cmd{label} بیرون از محیط، برچسب
    تکراری، و ارجاع به عنصر قبل از تعریف
    است.
\end{itemize}

این فصل، آخرین فصل از بخش دوم (اجزای کتاب)
بود. در بخش سوم (شخصی‌سازی)، یاد می‌گیرید
که چگونه \lr{MatinBook} را برای نیازهای
خود تنظیم کنید.